% FILE: 00_project_identity.tex
% Project: experiments-visualizer-nextjs
% Role: web-based-pi-visualization-surface

\section{Project Identity}
\label{sec:experiments-visualizer-nextjs-identity}

\subsection{Purpose}
\label{sec:experiments-visualizer-nextjs-identity-purpose}

The \texttt{experiments/visualizer-nextjs} project is the intended web-based process intelligence
visualization surface for the research program housed in \texttt{{\textasciitilde}/process-intelligence}.
Its purpose is to render, in a browser, the full range of process-intelligence evidence that the
research program manufactures: Petri net replay states, A{*} alignment move sequences, EWMA
drift-detection time series, and cryptographic receipt ledger blocks. The project is architecturally
positioned as the human-facing observation layer that sits above the Rust and Python computation
cores (\texttt{wasm4pm}, \texttt{pm4py}), translating machine-produced process evidence into
interactive visual artefacts inspectable by a researcher or auditor.

The project was bootstrapped via \texttt{create-next-app} targeting Next.js 16.2.6 and React 19.2.4,
establishing the scaffold for a TypeScript/React application. At the time of evidence extraction
the application had been compiled to a \texttt{.next} production build with
\texttt{BUILD\_ID: 9NFA-LnYiAijWLGnz8L79}, confirming that the underlying framework integrates
without error. However, the application's primary page (\texttt{src/app/page.tsx}) remains the
stock \texttt{create-next-app} welcome page; no process-intelligence-specific React components
have been implemented.

\subsection{Repository Surface}
\label{sec:experiments-visualizer-nextjs-identity-surface}

The project resides at \texttt{/Users/sac/process-intelligence/experiments/visualizer-nextjs/}
inside the broader \texttt{experiments/} corpus. The directory structure follows the Next.js 16
App Router convention: page and layout components live under \texttt{src/app/}, global styles
are declared in \texttt{src/app/globals.css}, and build output is emitted to \texttt{.next/}.
Turbopack is the configured bundler, confirmed by the \texttt{build-diagnostics.json} record
(\texttt{buildStage: static-generation}, \texttt{useBuildWorker: true}).

The most substantive source artefact is \texttt{globals.css} at 973 lines. This file encodes a
complete dark-mode dashboard design system using CSS custom properties and class-based component
tokens. The design system covers five distinct visual subsystems: a Petri net SVG panel, an A{*}
process alignment block sequence, an EWMA drift-detection chart with alarm banners, a
cryptographic ledger explorer with tamper-detection UI states, and a primary dashboard grid
layout with a 320-pixel sidebar and a fluid main content area.

\subsection{Primary Research Role}
\label{sec:experiments-visualizer-nextjs-identity-role}

Within the thesis the project occupies the role of \textbf{web-based-pi-visualization-surface}.
Process intelligence research produces evidence in machine-readable form --- OCEL event logs,
conformance fitness scores, receipt JSON blobs, alignment traces. The visualizer-nextjs project
is designed to be the surface through which that evidence becomes human-auditable. In the full
architecture the visualizer is the terminal layer of the evidence chain: Raw event data is
admitted to a process model, replayed against it, scored for fitness, and the resulting
provenance chain is written to a cryptographic receipt ledger; the visualizer presents all of
these artefacts simultaneously in a single browser session.

The project also demonstrates that TypeScript and React can serve as a proof-of-concept UI
manufacturing substrate for process intelligence tooling, complementing the Rust execution core
and the Python oracle layer.

\subsection{Key Artifacts}
\label{sec:experiments-visualizer-nextjs-identity-artifacts}

\begin{itemize}
  \item \texttt{src/app/globals.css} (973 lines) --- the primary substantive artefact;
        encodes the complete visual vocabulary for all five process-intelligence UI subsystems.
  \item \texttt{.next/BUILD\_ID} (value: \texttt{9NFA-LnYiAijWLGnz8L79}) --- confirms
        successful compilation of the Next.js 16 scaffold against Turbopack.
  \item \texttt{.next/diagnostics/build-diagnostics.json} --- records
        \texttt{buildStage: static-generation} and \texttt{useBuildWorker: true}.
  \item \texttt{.next/server/app/index.html} --- the rendered static output of the current
        (placeholder) \texttt{page.tsx}.
  \item \texttt{src/app/page.tsx} --- the stock \texttt{create-next-app} welcome page;
        the implementation gap for the visualizer surface.
  \item \texttt{package.json} --- declares Next.js 16.2.6, React 19.2.4, TypeScript 5,
        Tailwind CSS v4 with \texttt{@tailwindcss/postcss}, and Geist/Geist Mono fonts.
  \item \texttt{AGENTS.md} and \texttt{CLAUDE.md} --- authoring-agent instructions noting
        breaking changes in Next.js 16 relative to prior training data.
\end{itemize}
